%Input preamble
%Style
\documentclass[12pt]{article}
\usepackage[top=1in, bottom=1in, left=1in, right=1in]{geometry}
\parindent 22pt
\usepackage{fancyhdr}

%Packages
\usepackage{adjustbox}
\usepackage{amsmath}
\usepackage{amsfonts}
\usepackage{amssymb}
\usepackage{bm}
\usepackage[table]{xcolor}
\usepackage{tabu}
\usepackage{color,soul}
\usepackage{makecell}
\usepackage{longtable}
\usepackage{multirow}
\usepackage[normalem]{ulem}
\usepackage{etoolbox}
\usepackage{graphicx}
\usepackage{tabularx}
\usepackage{ragged2e}
\usepackage{booktabs}
\usepackage{caption}
\usepackage{fixltx2e}
\usepackage[para, flushleft]{threeparttablex}
\usepackage[capposition=top,objectset=centering]{floatrow}
\usepackage{subcaption}
\usepackage{pdfpages}
\usepackage{pdflscape}
\usepackage{natbib}
\usepackage{bibunits}
\definecolor{maroon}{HTML}{990012}
\usepackage[colorlinks=true,linkcolor=maroon,citecolor=maroon,urlcolor=maroon,anchorcolor=maroon]{hyperref}
\usepackage{marvosym}
\usepackage{makeidx}
\usepackage{tikz}
\usetikzlibrary{shapes}
\usepackage{setspace}
\usepackage{enumerate}
\usepackage{rotating}
\usepackage{tocloft}
\usepackage{epstopdf}
\usepackage[titletoc]{appendix}
\usepackage{framed}
\usepackage{comment}
\usepackage{xr}
\usepackage{titlesec}
\usepackage{footnote}
\usepackage{longtable}
\newlength{\tablewidth}
\setlength{\tablewidth}{9.3in}
\setcounter{secnumdepth}{4}

\titleformat{\paragraph}
{\normalfont\normalsize\bfseries}{\theparagraph}{1em}{}
\titlespacing*{\paragraph}
{0pt}{3.25ex plus 1ex minus .2ex}{1.5ex plus .2ex}
\makeatletter
\pretocmd\start@align
{%
  \let\everycr\CT@everycr
  \CT@start
}{}{}
\apptocmd{\endalign}{\CT@end}{}{}
\makeatother
%Watermark
\usepackage[printwatermark]{xwatermark}
\usepackage{lipsum}
\definecolor{lightgray}{RGB}{220,220,220}
%\newwatermark[allpages,color=lightgray,angle=45,scale=3,xpos=0,ypos=0]{Preliminary Draft}

%Further subsection level
\usepackage{titlesec}
\setcounter{secnumdepth}{4}
\titleformat{\paragraph}
{\normalfont\normalsize\bfseries}{\theparagraph}{1em}{}
\titlespacing*{\paragraph}
{0pt}{3.25ex plus 1ex minus .2ex}{1.5ex plus .2ex}

\setcounter{secnumdepth}{5}
\titleformat{\subparagraph}
{\normalfont\normalsize\bfseries}{\thesubparagraph}{1em}{}
\titlespacing*{\subparagraph}
{0pt}{3.25ex plus 1ex minus .2ex}{1.5ex plus .2ex}

%Functions
\DeclareMathOperator{\cov}{Cov}
\DeclareMathOperator{\corr}{Corr}
\DeclareMathOperator{\var}{Var}
\DeclareMathOperator{\plim}{plim}
\DeclareMathOperator*{\argmin}{arg\,min}
\DeclareMathOperator*{\argmax}{arg\,max}

%Math Environments
\newtheorem{theorem}{Theorem}
\newtheorem{claim}{Claim}
\newtheorem{condition}{Condition}
\renewcommand\thecondition{C--\arabic{condition}}
\newtheorem{algorithm}{Algorithm}
\newtheorem{assumption}{Assumption}
\renewcommand\theassumption{A--\arabic{assumption}}
\newtheorem{remark}{Remark}
\renewcommand\theremark{R--\arabic{remark}}
\newtheorem{definition}[theorem]{Definition}
\newtheorem{hypothesis}[theorem]{Hypothesis}
\newtheorem{property}[theorem]{Property}
\newtheorem{example}[theorem]{Example}
\newtheorem{result}[theorem]{Result}
\newenvironment{proof}{\textbf{Proof:}}{$\bullet$}

%Commands
\newcommand\independent{\protect\mathpalette{\protect\independenT}{\perp}}
\def\independenT#1#2{\mathrel{\rlap{$#1#2$}\mkern2mu{#1#2}}}
\newcommand{\overbar}[1]{\mkern 1.5mu\overline{\mkern-1.5mu#1\mkern-1.5mu}\mkern 1.5mu}
\newcommand{\equald}{\ensuremath{\overset{d}{=}}}
\captionsetup[table]{skip=10pt}
%\makeindex

\setlength\parindent{20pt}
\setlength{\parskip}{0pt}

\newcolumntype{L}[1]{>{\raggedright\let\newline\\\arraybackslash\hspace{0pt}}m{#1}}
\newcolumntype{C}[1]{>{\centering\let\newline\\\arraybackslash\hspace{0pt}}m{#1}}
\newcolumntype{R}[1]{>{\raggedleft\let\newline\\\arraybackslash\hspace{0pt}}m{#1}}



%Logo
%\AddToShipoutPictureBG{%
%  \AtPageUpperLeft{\raisebox{-\height}{\includegraphics[width=1.5cm]{uchicago.png}}}
%}

\newcolumntype{L}[1]{>{\raggedright\let\newline\\\arraybackslash\hspace{0pt}}m{#1}}
\newcolumntype{C}[1]{>{\centering\let\newline\\\arraybackslash\hspace{0pt}}m{#1}}
\newcolumntype{R}[1]{>{\raggedleft\let\newline\\\arraybackslash\hspace{0pt}}m{#1}}

\newcommand{\mr}{\multirow}
\newcommand{\mc}{\multicolumn}

%\newcommand{\comment}[1]{}


\begin{document}

\subsection{Using our Estimates to Understand Recent Cost-benefit Analyses}

\noindent An example of the approach recent studies take on cost-benefit analyses is \citet{Kline-Walters_2016_QJE}. They use data of the Head Start Impact Study (HSIS) to evaluate Head Start and find that its benefit/cost between $1.50$ and $1.84$.\footnote{HSIS is a one-year-long randomized version of Head Start.} They proceed in three steps: (i) calculate the treatment effect on Wechsler Preschool and Primary Scale of Intelligence (WPPSI) at age 5; (ii) monetize this gain using the return to WPPSI at age 5 in terms of net present value of earnings at age 27 \citep{Chetty_Friedman_etal_2011_QJoE}. Calculations from \citet{Chetty_Friedman_etal_2011_QJoE} indicate that a 1 standard deviation gain in WPPSI at age 5 implies a $13.1\%$ increase in the net present value of earnings at age 27;\footnote{This is based on combining information from Project Star and administrative data at age 27.} and (iii) calculate the benefit-to-cost ratio based on this gain and their own calculations of the program's cost. This calculation is based on assigning the net present value of earnings at age 27 of $\$385,907.17$ to the control-group participants, which is provided by  \citet{Chetty_Friedman_etal_2011_QJoE}.\footnote{All money values that we provide in this section are in 2014 USD.}\\


\begin{table}[H] 
\begin{threeparttable}
\caption{Alternative Cost-benefit Analyses Calculations}
\label{table:comparing}
\centering 
\footnotesize

\begin{tabular}{cllcc}
\toprule
Age & \mc{1}{c}{NPV Source} & Component & \citet{Kline_Walters_2016_QJE} & Authors' Method \\
& & & Method & \\
\midrule
\multirow{2}{*}{27} & \cite{Chetty_Friedman_etal_2011_QJoE} & Labor income & 0.58 (s.e. 0.28) &  \\
& ABC/CARE-calculated & Labor income & 0.09 (s.e. 0.04) &  1.09 (s.e. 0.04)\\
\midrule
\multirow{2}{*}{34} & ABC/CARE-calculated & Labor income & 0.37 (s.e. 0.04) & 0.15 (s.e. 0.05) \\
& ABC/CARE-calculated & All & 1.21 (s.e. 0.05) &  3.20 (s.e. 1.04) \\
\midrule
\multirow{2}{*}{Life-cycle} &  ABC/CARE-calculated & Labor income & 1.56 (s.e. 0.08) & 1.55 (s.e. 0.76) \\
& ABC/CARE-calculated & All & 3.80 (s.e. 0.29) & 7.33 (s.e. 1.84) \\
\bottomrule
\end{tabular}

\begin{tablenotes}
\footnotesize
\item Note: This table displays displays benefit-cost ratios based on the methodology \citet{Kline-Walters_2016_QJE} and based on our own methodology. Age: age at which we stop calculating the net-present value. NPV Source: source where we obtain the net present value. Component: item used to compute net present value (all refers to the net present value of all the components). \citet{Kline-Walters_2016_QJE}: estimate based on these authors methodology. Author's method: estimates based on our methodology.
\end{tablenotes}
\end{threeparttable}
\end{table}

\noindent To analyze how our estimates compares to the method in \citet{Kline-Walters_2016_QJE}, we present a series of exercises in the third column of Table~\ref{table:comparing}. For comparison, the fourth column of Table~\ref{table:comparing} displays the analogous exercise using our own method.\\

\noindent In the first exercise, we use both the ``return to WPPSI'' and the net-present value of earnings at age 27 in  \citet{Chetty_Friedman_etal_2011_QJoE}. In the second exercise, we perform a similar exercise but we use our own estimate of the net-present value of earnings at age 27.\footnote{This allows us to compute our own ``return to WPPSI'' and impute it to the treatment-group individuals.} The reminder of exercises are similar, but (i) vary the age up to where we consider the net-present value of earnings; (ii) consider the value of all the components we analyze throughout the paper; or (iii) both.\\

\noindent Our methodology has three gains: (i) provides a more precise estimate of the net-present value (and the return to WPPSI) of the components, as our interpolation and extrapolation is based on matching our experimental sample to various non-experimental samples. We do not impute a net present value or a return of an \textit{ad hoc} experiment; (ii) we better quantify the effects of the analyzed experiment by considering the whole life-cycle; (iii) we better approximate the statistical uncertainty of our estimates by considering both the sampling error in the experimental and non-experimental samples and the forecasting error due to the interpolation and extrapolation.

%References
\singlespace
\bibliographystyle{chicago}
\bibliography{heckman}
\end{document}